\documentclass[12pt]{article}
\usepackage{amsmath, amssymb, amsthm}
\usepackage{geometry}
\usepackage{hyperref}
\usepackage{xcolor}
\usepackage{tcolorbox}
\usepackage{enumitem}
\usepackage{mdframed}
\usepackage{array}

\geometry{a4paper, margin=2.5cm}

\definecolor{problemblue}{RGB}{30, 80, 160}
\definecolor{solngreen}{RGB}{20, 120, 60}
\definecolor{notegray}{RGB}{80, 80, 80}
\definecolor{boxbg}{RGB}{240, 245, 255}

\tcbuselibrary{skins, breakable}

\newtcolorbox{problembox}[1]{
  colback=boxbg, colframe=problemblue,
  fonttitle=\bfseries\color{white},
  title={#1},
  boxrule=1pt, arc=4pt
}

\newtcolorbox{answerbox}{
  colback=green!8, colframe=solngreen,
  fonttitle=\bfseries,
  title={Final Answer},
  boxrule=1pt, arc=4pt
}

\newtcolorbox{propbox}{
  colback=yellow!10, colframe=orange!70!black,
  fonttitle=\bfseries,
  title={Properties / Concepts Used},
  boxrule=0.8pt, arc=4pt
}

\title{\textbf{JEE Advanced — Matrices \& Determinants}}
\author{Virat Dixit}
\date{\today}

\begin{document}

\maketitle

% =====================================================================
\section{Problem 1 — JEE Advanced 2016}
% =====================================================================

\begin{problembox}{Problem 1 \hfill \normalfont\small [JEE Advanced 2016]}
Let $\alpha,\,\lambda,\,\mu \in \mathbb{R}$. Consider the system of linear equations:
\begin{align*}
\alpha x + 2y &= \lambda \\
3x - 2y &= \mu
\end{align*}
Which of the following statements is/are correct?

\begin{enumerate}[label=(\roman*)]
  \item If $\alpha = -3$, then the system has infinitely many solutions for all $\lambda$ and $\mu$.
  \item If $\alpha \neq -3$, then the system has a unique solution for all $\lambda$ and $\mu$.
  \item If $\lambda + \mu = 0$, then the system has infinitely many solutions for $\alpha = -3$.
  \item If $\lambda + \mu \neq 0$, then the system has no solution for $\alpha = -3$.
\end{enumerate}
\end{problembox}

\subsection*{Solution}

\textbf{Setting up the coefficient matrix.}

Let the coefficient matrix be
\[
A = \begin{pmatrix} \alpha & 2 \\ 3 & -2 \end{pmatrix}.
\]
Then
\[
\det(A) = \alpha(-2) - 3(2) = -2\alpha - 6 = -2(\alpha + 3).
\]

\subsubsection*{Checking Statement (ii) — Unique solution when $\alpha \neq -3$}

For a unique solution we need $\det(A) \neq 0$, i.e.,
\[
-2(\alpha+3) \neq 0 \implies \alpha \neq -3.
\]
Hence statement (ii) is \textbf{correct}.

\subsubsection*{Analysis when $\alpha = -3$}

When $\alpha = -3$, the system becomes:
\[
-3x + 2y = \lambda \quad \xrightarrow{\times(-1)} \quad 3x - 2y = -\lambda
\]
\[
3x - 2y = \mu
\]
Comparing the two equations, both reduce to $3x - 2y = -\lambda$ and $3x - 2y = \mu$.

\medskip
\noindent\textbf{Case 1: } $-\lambda = \mu$, i.e., $\lambda + \mu = 0$.\\
Both equations are identical $\Rightarrow$ the system is consistent with infinitely many solutions.\\
$\therefore$ Statement (iii) is \textbf{correct}.

\medskip
\noindent\textbf{Case 2: } $\lambda + \mu \neq 0$.\\
The two equations are contradictory (same LHS, different RHS) $\Rightarrow$ no solution.\\
$\therefore$ Statement (iv) is \textbf{correct}.

\subsubsection*{Checking Statement (i)}

Statement (i) claims infinitely many solutions for \emph{all} $\lambda, \mu$ when $\alpha = -3$. But from Case 2 above, if $\lambda + \mu \neq 0$ the system has \emph{no} solution. So statement (i) is \textbf{false}.

\begin{answerbox}
Correct statements: \textbf{(ii), (iii), (iv)}.
\end{answerbox}

\begin{propbox}
\begin{enumerate}[leftmargin=*, label=\textbf{P\arabic*.}]
  \item \textbf{Condition for unique solution:} $\det(A) \neq 0$.
  \item \textbf{Condition for no / infinitely many solutions:} $\det(A) = 0$; consistency then decides between the two cases.
  \item \textbf{Consistency check for singular systems:} Multiply one equation to match the other and compare RHS. Equal $\Rightarrow$ infinitely many solutions; unequal $\Rightarrow$ no solution.
  \item \textbf{Negation of a universal statement:} Statement (i) claims the result for \emph{all} $\lambda, \mu$; a single counterexample ($\lambda + \mu \neq 0$) suffices to disprove it.
\end{enumerate}
\end{propbox}

% =====================================================================
\section{Problem 2 — JEE Advanced 2017}
% =====================================================================

\begin{problembox}{Problem 2 \hfill \normalfont\small [JEE Advanced 2017]}
For a real number $\alpha$, if the system
\[
\begin{pmatrix} 1 & \alpha & \alpha^2 \\ \alpha & 1 & \alpha \\ \alpha^2 & \alpha & 1 \end{pmatrix}
\begin{pmatrix} x \\ y \\ z \end{pmatrix}
=
\begin{pmatrix} 1 \\ -1 \\ 1 \end{pmatrix}
\]
of linear equations has infinitely many solutions, then find $1 + \alpha + \alpha^2$.
\end{problembox}

\subsection*{Solution}

\textbf{Step 1: Condition for infinitely many solutions.}

For infinitely many solutions, the coefficient matrix must be singular \emph{and} the system must be consistent:
\[
\det(A) = 0 \quad \text{and the system is consistent.}
\]

\textbf{Step 2: Computing the determinant.}

\[
\Delta = \begin{vmatrix} 1 & \alpha & \alpha^2 \\ \alpha & 1 & \alpha \\ \alpha^2 & \alpha & 1 \end{vmatrix}
\]

Expanding along the first row:
\begin{align*}
\Delta &= 1\cdot(1 - \alpha^2) - \alpha(\alpha - \alpha^3) + \alpha^2(\alpha^2 - \alpha^2) \\
       &= (1 - \alpha^2) - \alpha^2(1 - \alpha^2) + \alpha^2(\alpha^2 - \alpha^2) \\
       &= (1-\alpha^2)(1 - \alpha^2) \\[4pt]
\Delta &= 1 - \alpha^2 - \alpha^2 + \alpha^4 = 1 - 2\alpha^2 + \alpha^4 = (1 - \alpha^2)^2.
\end{align*}

\textbf{Step 3: Setting $\Delta = 0$.}

\[
(1 - \alpha^2)^2 = 0 \implies \alpha^2 = 1 \implies \alpha = \pm 1.
\]

\textbf{Step 4: Consistency check for $\alpha = 1$.}

The matrix becomes all 1's:
\[
\begin{pmatrix}1&1&1\\1&1&1\\1&1&1\end{pmatrix}
\begin{pmatrix}x\\y\\z\end{pmatrix}
=\begin{pmatrix}1\\-1\\1\end{pmatrix}
\]
Equations: $x+y+z=1$, $x+y+z=-1$, $x+y+z=1$. This is \textbf{inconsistent} (contradictory RHS). Reject $\alpha = 1$.

\textbf{Step 5: Consistency check for $\alpha = -1$.}

The matrix becomes:
\[
\begin{pmatrix}1&-1&1\\-1&1&-1\\1&-1&1\end{pmatrix}
\begin{pmatrix}x\\y\\z\end{pmatrix}
=\begin{pmatrix}1\\-1\\1\end{pmatrix}
\]
All three equations reduce to $x - y + z = 1$. These are \textbf{dependent} $\Rightarrow$ infinitely many solutions. \checkmark

Hence $\alpha = -1$.

\textbf{Step 6: Finding the required expression.}

\[
1 + \alpha + \alpha^2 = 1 + (-1) + (-1)^2 = 1 - 1 + 1 = \boxed{1}
\]

\begin{answerbox}
\[1 + \alpha + \alpha^2 = \mathbf{1}\]
\end{answerbox}

\begin{propbox}
\begin{enumerate}[leftmargin=*, label=\textbf{P\arabic*.}]
  \item \textbf{Singular matrix condition:} $\det(A) = 0$ is necessary but not sufficient for infinitely many solutions; consistency must be separately verified.
  \item \textbf{$3\times 3$ determinant expansion:} Along the first row using cofactors.
  \item \textbf{Factoring:} $1 - 2\alpha^2 + \alpha^4 = (1 - \alpha^2)^2$ — recognising a perfect square saves computation.
  \item \textbf{Consistency via row reduction:} When $\det = 0$, check whether the augmented matrix leads to a contradiction ($0 = c \neq 0$) or dependence ($0 = 0$).
\end{enumerate}
\end{propbox}

% =====================================================================
\section{Problem 3 — JEE Advanced 2018}
% =====================================================================

\begin{problembox}{Problem 3 \hfill \normalfont\small [JEE Advanced 2018]}
Let $P$ be a $3\times 3$ matrix such that all entries of $P$ are from the set $\{-1,\,0,\,1\}$. Find the maximum possible value of $\det(P)$.
\end{problembox}

\subsection*{Solution}

\textbf{Step 1: Determinant expansion idea.}

For a $3\times 3$ matrix, the determinant is a sum of 6 signed products (by Sarrus or Leibniz):
\[
\det A = a_{11}a_{22}a_{33} + a_{12}a_{23}a_{31} + a_{13}a_{21}a_{32}
       - a_{13}a_{22}a_{31} - a_{12}a_{21}a_{33} - a_{11}a_{23}a_{32}.
\]
Each of the 6 products involves entries from $\{-1,0,1\}$, so each product is at most $1$ in magnitude. Thus naively $|\det A| \leq 6$. However, we use a sharper bound.

\textbf{Step 2: Hadamard's Inequality.}

If the rows of $A$ are $R_1, R_2, R_3 \in \mathbb{R}^3$, then
\[
|\det A| \leq \|R_1\|\,\|R_2\|\,\|R_3\|.
\]
Since each entry is from $\{-1,0,1\}$, the maximum Euclidean norm of any row is achieved when all three entries are $\pm 1$:
\[
\|R_i\| \leq \sqrt{1^2 + 1^2 + 1^2} = \sqrt{3}.
\]
Therefore
\[
|\det A| \leq (\sqrt{3})^3 = 3\sqrt{3} \approx 5.196.
\]
Since $\det A$ is always an integer (integer entries), we get
\[
|\det A| \leq 5.
\]

\textbf{Step 3: Is $\det A = 5$ achievable?}

Hadamard's bound is tight only when rows are mutually orthogonal. For $\pm 1$ entries in dimension 3, no three mutually orthogonal vectors exist (this is related to the non-existence of a $3\times 3$ Hadamard matrix with $\pm 1$ entries). Hence $|\det A| = 5$ is \textbf{not achievable}.

\textbf{Step 4: Try $\det A = 4$.}

We need to find an explicit matrix achieving $|\det| = 4$. Consider:
\[
A = \begin{pmatrix} 1 & 1 & 0 \\ 1 & -1 & 1 \\ 1 & 0 & -1 \end{pmatrix}.
\]
Expanding along the first row:
\begin{align*}
\det A &= 1\cdot\bigl((-1)(-1)-(1)(0)\bigr) - 1\cdot\bigl((1)(-1)-(1)(1)\bigr) + 0 \\
       &= 1(1) - 1(-2) = 1 + 2 = 3.
\end{align*}

We can do better. Consider instead:
\[
A = \begin{pmatrix} 1 & 1 & 0 \\ 1 & -1 & 1 \\ 1 & 0 & -1 \end{pmatrix}
\xrightarrow{\text{try}} 
A = \begin{pmatrix} 1 & 1 & 0 \\ 1 & -1 & 1 \\ 0 & 1 & -1 \end{pmatrix}.
\]

The matrix from the notebook achieving $\det = 4$:
\[
A = \begin{pmatrix} 1 & 1 & 0 \\ 1 & -1 & 1 \\ 1 & 0 & -1 \end{pmatrix}.
\]

A clean verified example: take
\[
A = \begin{pmatrix} 1 & 1 & 1 \\ 1 & -1 & 1 \\ 1 & 1 & -1 \end{pmatrix}.
\]
\begin{align*}
\det A &= 1\bigl((-1)(-1)-(1)(1)\bigr) - 1\bigl((1)(-1)-(1)(1)\bigr) + 1\bigl((1)(1)-(-1)(1)\bigr)\\
       &= 1(1-1) - 1(-1-1) + 1(1+1)\\
       &= 0 + 2 + 2 = 4.
\end{align*}

So $\det = 4$ is \textbf{achievable}. Since 5 is ruled out, the maximum is 4.

\begin{answerbox}
The maximum possible value of $\det(P)$ is $\mathbf{4}$, achieved for example by:
\[
P = \begin{pmatrix} 1 & 1 & 1 \\ 1 & -1 & 1 \\ 1 & 1 & -1 \end{pmatrix}, \quad \det P = 0 + 2 + 2 = 4.
\]
\end{answerbox}

\begin{propbox}
\begin{enumerate}[leftmargin=*, label=\textbf{P\arabic*.}]
  \item \textbf{Leibniz / Sarrus formula:} For a $3\times 3$ matrix, $\det$ is a signed sum of 6 products of entries.
  \item \textbf{Hadamard's Inequality:} $|\det A| \leq \|R_1\|\|R_2\|\|R_3\|$. Gives an upper bound of $3\sqrt{3} \approx 5.196$, which forces $|\det A| \leq 5$ for integer determinants.
  \item \textbf{Integer determinant:} Since entries are integers, $\det A \in \mathbb{Z}$, so the Hadamard bound $3\sqrt{3}$ immediately reduces to $|\det A| \leq 5$.
  \item \textbf{Non-existence of $3\times 3$ Hadamard matrix:} Three mutually orthogonal $\pm 1$ vectors cannot exist in $\mathbb{R}^3$, which rules out $|\det A| = 5$.
  \item \textbf{Construction to confirm attainability:} After ruling out 5, an explicit matrix with $\det = 4$ confirms the maximum.
\end{enumerate}
\end{propbox}

% =====================================================================
\section*{Summary of Key Matrix / Determinant Properties}
% =====================================================================

\begin{mdframed}[backgroundcolor=boxbg, linecolor=problemblue, linewidth=1pt]
\textbf{Properties Consolidation}
\begin{align*}
&\bullet\ \det(A) \neq 0 \Leftrightarrow \text{unique solution} && \text{(invertibility criterion)}\\
&\bullet\ \det(A) = 0,\ \text{consistent} \Leftrightarrow \text{infinitely many solutions} && \text{(singular consistent)}\\
&\bullet\ \det(A) = 0,\ \text{inconsistent} \Leftrightarrow \text{no solution} && \text{(singular inconsistent)}\\
&\bullet\ |\det A| \leq \|R_1\|\|R_2\|\|R_3\| && \text{(Hadamard's inequality)}\\
&\bullet\ \det A \in \mathbb{Z}\ \text{when entries} \in \mathbb{Z} && \text{(integer determinant)}
\end{align*}
\end{mdframed}

\end{document}
